\documentclass[ twoside,openright,titlepage,numbers=noenddot,headinclude,%1headlines,% letterpaper a4paper
                footinclude=true,cleardoublepage=empty,abstractoff, % <--- obsolete, remove (todo)
                BCOR=5mm,paper=a4,fontsize=11pt,%11pt,a4paper,%
                ngerman,american,%
                ]{scrreprt}


%load fonts und useful packages. etc.
% ****************************************************************************************************
% classicthesis-config.tex 
% formerly known as loadpackages.sty, classicthesis-ldpkg.sty, and classicthesis-preamble.sty 
% Use it at the beginning of your ClassicThesis.tex, or as a LaTeX Preamble 
% in your ClassicThesis.{tex,lyx} with \input{classicthesis-config}
% ****************************************************************************************************  
% If you like the classicthesis, then I would appreciate a postcard. 
% My address can be found in the file ClassicThesis.pdf. A collection 
% of the postcards I received so far is available online at 
% http://postcards.miede.de
% ****************************************************************************************************

% ****************************************************************************************************
% 1. Configure classicthesis for your needs here, e.g., remove "drafting" below 
% in order to deactivate the time-stamp on the pages
% ****************************************************************************************************
\PassOptionsToPackage{eulerchapternumbers,listings,drafting,%
				 pdfspacing,%floatperchapter,%linedheaders,%
				 subfig,beramono,eulermath}{classicthesis}										
% ********************************************************************
% Available options for classicthesis.sty 
% (see ClassicThesis.pdf for more information):
% drafting
% parts nochapters linedheaders
% eulerchapternumbers beramono eulermath pdfspacing minionprospacing
% tocaligned dottedtoc manychapters
% listings floatperchapter subfig
% ********************************************************************

% ********************************************************************
% Triggers for this config
% ******************************************************************** 
\usepackage{ifthen}
\newboolean{enable-backrefs} % enable backrefs in the bibliography
\setboolean{enable-backrefs}{false} % true false
% ****************************************************************************************************


% ****************************************************************************************************
% 2. Personal data and user ad-hoc commands
% ****************************************************************************************************
\newcommand{\myTitle}{Flight Path Optimisation\xspace}
\newcommand{\mySubtitle}{ Project Report\xspace}
\newcommand{\myDegree}{\xspace}
\newcommand{\myName}{Bob Vergauwen and Moritz Wolter\xspace}
\newcommand{\myProf}{Prof. dr. ir. Toon van Waterschoot \xspace}
\newcommand{\myOtherProf}{\xspace}
\newcommand{\mySupervisor}{\xspace}
\newcommand{\myFaculty}{\xspace}
\newcommand{\myDepartment}{\xspace}
\newcommand{\myUni}{KU Leuven\xspace}
\newcommand{\myLocation}{Leuven\xspace}
\newcommand{\myTime}{	\selectlanguage{american}
			\today \xspace}
\newcommand{\myVersion}{Version 1\xspace}

% ********************************************************************
% Setup, finetuning, and useful commands
% ********************************************************************
\newcounter{dummy} % necessary for correct hyperlinks (to index, bib, etc.)
\newlength{\abcd} % for ab..z string length calculation
\providecommand{\mLyX}{L\kern-.1667em\lower.25em\hbox{Y}\kern-.125emX\@}
\newcommand{\ie}{i.\,e.}
\newcommand{\Ie}{I.\,e.}
\newcommand{\eg}{e.\,g.}
\newcommand{\Eg}{E.\,g.} 
% ****************************************************************************************************


% ****************************************************************************************************
% 3. Loading some handy packages
% ****************************************************************************************************
% ******************************************************************** 
% Packages with options that might require adjustments
% ******************************************************************** 
\PassOptionsToPackage{utf8}{inputenc}	% latin9 (ISO-8859-9) = latin1+"Euro sign"
 \usepackage{inputenc}				

\PassOptionsToPackage{english}{babel}   % change this to your language(s)
% Spanish languages need extra options in order to work with this template
%\PassOptionsToPackage{spanish,es-lcroman}{babel}
 \usepackage{babel}					

 %\PassOptionsToPackage{square,numbers}{natbib}
 %\usepackage{natbib}
 \usepackage[fixlanguage]{babelbib}
 \selectbiblanguage{german}
 \bibliographystyle{babplain}

\PassOptionsToPackage{fleqn}{amsmath}		% math environments and more by the AMS 
 \usepackage{amsmath}

%When using tikZ one must also use the color package.
\PassOptionsToPackage{usenames,dvipsnames}{color}
 \usepackage{color}

%This package lets you compile tikz graphics.
 \usepackage{tikz}
 \usepackage{pgfplots}
 \usetikzlibrary{decorations.markings}
 \usepackage{standalone}

% ******************************************************************** 
% General useful packages
% ******************************************************************** 
\PassOptionsToPackage{T1}{fontenc} % T2A for cyrillics
	\usepackage{fontenc}     
\usepackage{textcomp} % fix warning with missing font shapes
\usepackage{scrhack} % fix warnings when using KOMA with listings package          
\usepackage{xspace} % to get the spacing after macros right  
\usepackage{mparhack} % get marginpar right
\usepackage{fixltx2e} % fixes some LaTeX stuff 
\PassOptionsToPackage{printonlyused,smaller}{acronym}
	\usepackage{acronym} % nice macros for handling all acronyms in the thesis
%\renewcommand*{\acsfont}[1]{\textssc{#1}} % for MinionPro
\renewcommand{\bflabel}[1]{{#1}\hfill} % fix the list of acronyms
% ****************************************************************************************************


% ****************************************************************************************************
% 4. Setup floats: tables, (sub)figures, and captions
% ****************************************************************************************************
\usepackage{tabularx} % better tables
	\setlength{\extrarowheight}{3pt} % increase table row height
\newcommand{\tableheadline}[1]{\multicolumn{1}{c}{\spacedlowsmallcaps{#1}}}
\newcommand{\myfloatalign}{\centering} % to be used with each float for alignment
\usepackage{caption}
\captionsetup{format=hang,font=small}
\usepackage{subfig}  
% ****************************************************************************************************


% ****************************************************************************************************
% 5. Setup code listings
% ****************************************************************************************************
\usepackage{listings} 
%\lstset{emph={trueIndex,root},emphstyle=\color{BlueViolet}}%\underbar} % for special keywords
\lstset{language=[LaTeX]Tex,%C++,
    keywordstyle=\color{RoyalBlue},%\bfseries,
    basicstyle=\small\ttfamily,
    %identifierstyle=\color{NavyBlue},
    commentstyle=\color{Green}\ttfamily,
    stringstyle=\rmfamily,
    numbers=none,%left,%
    numberstyle=\scriptsize,%\tiny
    stepnumber=5,
    numbersep=8pt,
    showstringspaces=false,
    breaklines=true,
    frameround=ftff,
    frame=single,
    belowcaptionskip=.75\baselineskip
    %frame=L
} 
% ****************************************************************************************************    		   


% ****************************************************************************************************
% 6. PDFLaTeX, hyperreferences and citation backreferences
% ****************************************************************************************************
% ********************************************************************
% Using PDFLaTeX
% ********************************************************************
\PassOptionsToPackage{pdftex,hyperfootnotes=false,pdfpagelabels}{hyperref}
	\usepackage{hyperref}  % backref linktocpage pagebackref
\pdfcompresslevel=9
\pdfadjustspacing=1 
\PassOptionsToPackage{pdftex}{graphicx}
	\usepackage{graphicx} 

% ********************************************************************
% Setup the style of the backrefs from the bibliography
% (translate the options to any language you use)
% ********************************************************************
\newcommand{\backrefnotcitedstring}{\relax}%(Not cited.)
\newcommand{\backrefcitedsinglestring}[1]{(Cited on page~#1.)}
\newcommand{\backrefcitedmultistring}[1]{(Cited on pages~#1.)}
\ifthenelse{\boolean{enable-backrefs}}%
{%
		\PassOptionsToPackage{hyperpageref}{backref}
		\usepackage{backref} % to be loaded after hyperref package 
		   \renewcommand{\backreftwosep}{ and~} % separate 2 pages
		   \renewcommand{\backreflastsep}{, and~} % separate last of longer list
		   \renewcommand*{\backref}[1]{}  % disable standard
		   \renewcommand*{\backrefalt}[4]{% detailed backref
		      \ifcase #1 %
		         \backrefnotcitedstring%
		      \or%
		         \backrefcitedsinglestring{#2}%
		      \else%
		         \backrefcitedmultistring{#2}%
		      \fi}%
}{\relax}    

% ********************************************************************
% Hyperreferences
% ********************************************************************
\hypersetup{%
    %draft,	% = no hyperlinking at all (useful in b/w printouts)
    colorlinks=true, linktocpage=true, pdfstartpage=3, pdfstartview=FitV,%
    % uncomment the following line if you want to have black links (e.g., for printing)
    %colorlinks=false, linktocpage=false, pdfborder={0 0 0}, pdfstartpage=3, pdfstartview=FitV,% 
    breaklinks=true, pdfpagemode=UseNone, pageanchor=true, pdfpagemode=UseOutlines,%
    plainpages=false, bookmarksnumbered, bookmarksopen=true, bookmarksopenlevel=1,%
    hypertexnames=true, pdfhighlight=/O,%nesting=true,%frenchlinks,%
    urlcolor=webbrown, linkcolor=RoyalBlue, citecolor=webgreen, %pagecolor=RoyalBlue,%
    %urlcolor=Black, linkcolor=Black, citecolor=Black, %pagecolor=Black,%
    pdftitle={\myTitle},%
    pdfauthor={\textcopyright\ \myName, \myUni, \myFaculty},%
    pdfsubject={},%
    pdfkeywords={},%
    pdfcreator={pdfLaTeX},%
    pdfproducer={LaTeX with hyperref and classicthesis}%
}   

% ********************************************************************
% Setup autoreferences
% ********************************************************************
% There are some issues regarding autorefnames
% http://www.ureader.de/msg/136221647.aspx
% http://www.tex.ac.uk/cgi-bin/texfaq2html?label=latexwords
% you have to redefine the makros for the 
% language you use, e.g., american, ngerman
% (as chosen when loading babel/AtBeginDocument)
% ********************************************************************
\makeatletter
\@ifpackageloaded{babel}%
    {%
       \addto\extrasamerican{%
					\renewcommand*{\figureautorefname}{Figure}%
					\renewcommand*{\tableautorefname}{Table}%
					\renewcommand*{\partautorefname}{Part}%
					\renewcommand*{\chapterautorefname}{Chapter}%
					\renewcommand*{\sectionautorefname}{Section}%
					\renewcommand*{\subsectionautorefname}{Section}%
					\renewcommand*{\subsubsectionautorefname}{Section}% 	
				}%
       \addto\extrasngerman{% 
					\renewcommand*{\paragraphautorefname}{Absatz}%
					\renewcommand*{\subparagraphautorefname}{Unterabsatz}%
					\renewcommand*{\footnoteautorefname}{Fu\"snote}%
					\renewcommand*{\FancyVerbLineautorefname}{Zeile}%
					\renewcommand*{\theoremautorefname}{Theorem}%
					\renewcommand*{\appendixautorefname}{Anhang}%
					\renewcommand*{\equationautorefname}{Gleichung}%        
					\renewcommand*{\itemautorefname}{Punkt}%
				}%	
			% Fix to getting autorefs for subfigures right (thanks to Belinda Vogt for changing the definition)
			\providecommand{\subfigureautorefname}{\figureautorefname}%  			
    }{\relax}
\makeatother


% ****************************************************************************************************
% 7. Last calls before the bar closes
% ****************************************************************************************************
% ********************************************************************
% Development Stuff
% ********************************************************************
\listfiles
%\PassOptionsToPackage{l2tabu,orthodox,abort}{nag}
%	\usepackage{nag}
%\PassOptionsToPackage{warning, all}{onlyamsmath}
%	\usepackage{onlyamsmath}

% ********************************************************************
% Last, but not least... 
% ********************************************************************
\usepackage{classicthesis} 
% Die guten TU-Farben...
\renewcommand\cftchapfont{\normalfont\color{TealBlue}}
\renewcommand\cftchappagefont{\normalfont\color{TealBlue}}

% ****************************************************************************************************


% ****************************************************************************************************
% 8. Further adjustments (experimental)
% ****************************************************************************************************
% ********************************************************************
% Changing the text area
% ********************************************************************
%\linespread{1.05} % a bit more for Palatino
%\areaset[current]{312pt}{761pt} % 686 (factor 2.2) + 33 head + 42 head \the\footskip
%\setlength{\marginparwidth}{7em}%
%\setlength{\marginparsep}{2em}%

% ********************************************************************
% Using different fonts
% ********************************************************************
%\usepackage[oldstylenums]{kpfonts} % oldstyle notextcomp
%\usepackage[osf]{libertine}
%\usepackage{hfoldsty} % Computer Modern with osf
%\usepackage[light,condensed,math]{iwona}
%\renewcommand{\sfdefault}{iwona}
%\usepackage{lmodern} % <-- no osf support :-(
%\usepackage[urw-garamond]{mathdesign} <-- no osf support :-(
% ****************************************************************************************************

\begin{document}
\frenchspacing
\raggedbottom
\selectlanguage{american} % american ngerman
%\renewcommand*{\bibname}{new name}
%\setbibpreamble{}
\pagenumbering{roman}
\pagestyle{plain}
%create titelpage
%*******************************************************
% Titlepage
%*******************************************************
\begin{titlepage}
	% if you want the titlepage to be centered, uncomment and fine-tune the line below (KOMA classes environment)
	\begin{addmargin}[-1cm]{-3cm}
    \begin{center}
        \large  

        \hfill

        \vfill

        \begingroup
            \color{TealBlue}\spacedallcaps{\myTitle} \\ \bigskip
        \endgroup

        \spacedlowsmallcaps{\myName}

        \vfill

        \includegraphics[width=6cm]{FrontBackmatter/logo_kuleuven.png} \\ \medskip

        \mySubtitle \\ \medskip   
        %\myDegree \\
        %\myDepartment \\                            
        %\myFaculty \\
	Supervised by \myProf \\
	\myOtherProf
        %\myUni \\ \bigskip

        \myTime\ -- \myVersion

        \vfill                      

    \end{center}  
  \end{addmargin}       
\end{titlepage}   

%%*******************************************************
% Table of Contents
%*******************************************************
%\phantomsection
\refstepcounter{dummy}
\pdfbookmark[1]{\contentsname}{tableofcontents}
\setcounter{tocdepth}{2} % <-- 2 includes up to subsections in the ToC
\setcounter{secnumdepth}{3} % <-- 3 numbers up to subsubsections
\manualmark
\markboth{\spacedlowsmallcaps{\contentsname}}{\spacedlowsmallcaps{\contentsname}}
\tableofcontents 
\automark[section]{chapter}
\renewcommand{\chaptermark}[1]{\markboth{\spacedlowsmallcaps{#1}}{\spacedlowsmallcaps{#1}}}
\renewcommand{\sectionmark}[1]{\markright{\thesection\enspace\spacedlowsmallcaps{#1}}}

\vspace*{8ex}
\noindent
Our sorce code my be downloaded from: \\
\url{https://github.com/v0lta/odeProject}

%*******************************************************
% List of Figures and of the Tables
%*******************************************************
\clearpage

\begingroup 
    \let\clearpage\relax
    \let\cleardoublepage\relax
    \let\cleardoublepage\relax
    %*******************************************************
    % List of Figures
    %*******************************************************    
    %\phantomsection 
    \refstepcounter{dummy}
    %\addcontentsline{toc}{chapter}{\listfigurename}
    \pdfbookmark[1]{\listfigurename}{lof}
    \listoffigures

    \vspace*{8ex}

    %*******************************************************
    % List of Tables
    %*******************************************************
    %\phantomsection 
    \refstepcounter{dummy}
    %\addcontentsline{toc}{chapter}{\listtablename}
    \pdfbookmark[1]{\listtablename}{lot}
    \listoftables
        
    \vspace*{8ex}
%   \newpage
    
    %*******************************************************
    % List of Listings
    %*******************************************************      
%	  %\phantomsection 
 %   \refstepcounter{dummy}
    %\addcontentsline{toc}{chapter}{\lstlistlistingname}
%    \pdfbookmark[1]{\lstlistlistingname}{lol}
%    \lstlistoflistings 

%    \vspace*{8ex}
       
    %*******************************************************
    % Acronyms
    %*******************************************************
    %\phantomsection 
    %\refstepcounter{dummy}
    %\pdfbookmark[1]{Acronyms}{acronyms}
    %\markboth{\spacedlowsmallcaps{Acronyms}}{\spacedlowsmallcaps{Acronyms}}
    %\chapter*{Acronyms}
    %\begin{acronym}[UML]
%	\acro{parfor}{parallel for Schleife}         
%	\acro{svd}{Singulärwertzerlegung}%
%	\acro{Matlab}{matrix laboratory}%
%	\acro{SISO}{Single-Input Single-Output}
%    \end{acronym}                     
\endgroup

\cleardoublepage

\pagestyle{scrheadings}


\pagenumbering{arabic}

\chapter{Motivation and Implementation}
\section{Introduction}
In this project we are going to find the optimal path for a solar powered airplane \footnote{Like the one from \url{http://www.solarimpulse.com/}}. During flight the plane's electrical engines drain power from it's battery, while at the same time solar panels on the plane's wings convert energy from the sun's rays into electric energy. However clouds may block the sun's rays from reaching the panels. In addition the sun's intensity increases when flying closer to the equator. In order to keep as much energy in the batteries as possible the flight path has to be optimized.

\section{Problem formulation and Analysis}

In this project two different models have been used, a time independent one and time depended model.
The big difference between the two models is the dimension of the parameter input space.
In the time invariant model this space consists of a $2 \times m $ matrix, this vector stands for the $ X $ and $ Y $ coordinates of the plane throughout the flight.
The time for one step is set to be constant.
This means that when the step size taken by the plane increases the velocity of the plane goes up as well.

To expand upon this model we introduced a third column vector as an input, a time vector.
This vector holds all of the time values associated with each step of the plane.
In this way can the path length be varied for each step without the need to accelerate.  


\subsection{Time invariant case}
In the optimization process we use randomly generated weather data and sun data corresponding to a flight between $48^\circ$ and $50^\circ$ latitude. We generate the weather data by cubicly interpolating a random matrix of size \texttt{n} over a grid. The larger n becomes the more complex the weather will be. This functionality is implemented in \texttt{Static\_Weather.m} \textbf{repo is still private...} \footnote{All matlab code can be found on \url{TODO: Put public cloned repo link here!}.} The class returns a weather data matrix $W$.

Data for the sun comes from a couple of established equations that relate the local time and position to the sun's intensity  which happens in \texttt{sun.m}. The class returns a sun data Matrix $S$. We define low values in the weather Data Matrix as cloudy and high values as sunny weather. As clouds reduce the electric yield of our plane's solar cells we define the overall intensity as:
\begin{equation}
$I$ \:= \: $W$ .* $S$
\end{equation}
While following a convention where $.*$ denotes element wise multiplication we end up with an intensity matrix $I$, which contains sun intensity data weighted according to the random weather.

Next we took a closer look at the airplane. Here we consider the degree to which the plane can recharge its battery during flight, which is the path integral over the intensity data:
\begin{equation} \label{eqn:sunNotime}
E(x) = \oint_{\Gamma} I  \; dx.
\end{equation}
With $dx$ being the length of each step. Having found the additional energy we are able to convert during flight, we now look at losses, which are due to drag and acceleration. We define acceleration cost in the time invariant case as the scalar product of the vectors containing local accelerations:
\begin{equation}
A(\mathbf{x}) = \mathbf{a}^T \cdot \mathbf{a}
\end{equation}
\marginpar{To calculate the derivatives in the time invariant case we approximate $dt$ as the total time over distance traveled.} However we only count cases where the acceleration is positive since deceleration does not drain battery power.
Finally we consider the drag force which the plain experiences on its path, which we approximate as the dot product of the local speed vector:
\begin{equation}
D(\mathbf{x}) = \mathbf{s}^T \cdot \mathbf{s}
\end{equation}
Having investigated the approximate battery power lost and gained on a specific path, we can now proceed to formulating a cost function for the time invariant case.

\section{Cost function formulation}
\subsection{Time invariant case}
In the time invariant case we define the optimal path as the path from start to destination where we loose as little battery power as possible. The problem parameter is a path vector $\mathbf{x} \in \mathbb{R}^m$. Which contains $m$ points on the path. Introducing weights for the different power gains and losses we define the cost function:
\begin{equation}
Energy(\mathbf{x}) = E(\mathbf{x}) -0.01A(\mathbf{x}) - 200D(\mathbf{x}).
\end{equation}
\begin{align}
       \min_{x,y}( -\text{Energy}(x,y,t)) & \hspace{5mm}  \text{s. t. } \\
		& x \in [0,1]^m\\
		& y \in [0,1]^m\\
		& [x[0],y[0]] = [x_0,y_0] \\
		& [x[m],y[m]] = [x_m,y_m] \\
\end{align}
Where the first constraint is that the path my not leave the weather data grid. And the second one that the path has to start from the place of origin and end at the destination. We solved this problem starting with the straight line from (0.1,0.5) to (0.9,0.5) as initial input using matlab's \texttt{fmincon} solver. This cost function is implemented in the \texttt{Airplane.m} class file.


\section{Time varying case}

In the previous section we performed an elaborate explanation of the a first model of the air plane.
This model did not include time  and it was assumed that the air plane flew with a constant speed.
This however is not realistic, we want to stay longer in sunny regions and fly fast trough cloudy ones.
To compensate for this need it is a good idea to include the time as an input parameter as well.
This new freedom will introduce a new dimension and thus make the optimization problem more complex.

\subsection{Formulation of the problem}

Lets start by reformulating the optimization problem in a formal way.

\begin{align}
       \min_{x,y,t }( -\text{Energy}(x,y,t)) & \hspace{5mm}  \text{s. t. } \\
		& x \in [0,1]^m\\
		& y \in [0,1]^m\\
		& [x[0],y[0]] = [x_0,y_0] \\
		& [x[m],y[m]] = [x_m,y_m] \\
		& t_0 = 0\\
		& t(i) \leq t(i+1)| \forall i \in \left\{1,2 \dots m\right\}
\end{align}
The conditions for the spacial coordinates have stayed the same, the plane has to stay in a certain region.
By introducing time dependency we have to make sure that the time is increasing throughout the flight, without this constraint the solution would be nonsense.

Lets now look at how the introduction of time changes the Energy function, we are still taking in to account the same quantities as in the previous model; the positive acceleration energy, the air resistance and the solar gain.
The only  thing that really changed is the way the speed of the air plane is defined.
In the previous model we assumed that each step taken by the plane took an equal time step regardless of the length of this step, this has changed in the time varying model, the time step is now a parameter that could be chosen optimal.
In a more mathematical formulation we took the speed to be defined as
\begin{equation}
v[i] = \dfrac{l[i+1]-l[i]}{t[i+1]-t[i]}
\end{equation}
in this formulation is $ l[i] $ the path length  from the beginning till the $i$-th point.

When this speed is used instead of the speed as defined in the above section most of the calculations stayed the same, except for the calculation of the total energy yield from the sun.
For the time independent model this calculations was done using Equation \ref{eqn:sunNotime}.
In this calculation we assumed that the velocity of the plane  constant was so that $ dx/ dt = v $.
This however is no longer true and the integration variable has to be adapted for it, so we get
\begin{equation} \label{eqn:sunNotime}
E(t) = \oint_{\Gamma} I  \; dt.
\end{equation}

The rest of the analysis stayed the same.


\subsection{Further expansion of the model}

A lot of assumptions has been made throughout the model.
Most of these are oversimplifications of the real world and we barley scratched the surface of the real optimisation problem behind the airplane.
Things we did not include in the model are for example a minimum speed and a maximum speed for the airplane.
We only calculated the residual energy at the end of the flight, but was the energy of the airplane at all times a positive value or did we use energy we did not have available?
Can we make the weather more realistic and make it time varying as well? If we did so than $ t[0] =0 $ was no longer a condition and we could start finding the ideal time of departure.

These are just a few possible modifications that could make the model more realistic and more usable.






\chapter{Results}
\section{Time invariant case}
\begin{figure}
% This file was created by matlab2tikz v0.4.7 running on MATLAB 8.4.
% Copyright (c) 2008--2014, Nico Schlömer <nico.schloemer@gmail.com>
% All rights reserved.
% Minimal pgfplots version: 1.3
% 
% The latest updates can be retrieved from
%   http://www.mathworks.com/matlabcentral/fileexchange/22022-matlab2tikz
% where you can also make suggestions and rate matlab2tikz.
% 
\documentclass[tikz]{standalone}
\usepackage{pgfplots}
\usepackage{grffile}
\pgfplotsset{compat=newest}
\usetikzlibrary{plotmarks}
\usepackage{amsmath}

\begin{document}
\begin{tikzpicture}

\begin{axis}[%
width=4.56842105263158in,
height=3.60315789473684in,
view={0}{90},
scale only axis,
every outer x axis line/.append style={white!15!black},
every x tick label/.append style={font=\color{white!15!black}},
xmin=0,
xmax=1,
xmajorgrids,
every outer y axis line/.append style={white!15!black},
every y tick label/.append style={font=\color{white!15!black}},
ymin=0,
ymax=1,
ymajorgrids,
every outer z axis line/.append style={white!15!black},
every z tick label/.append style={font=\color{white!15!black}},
zmin=-0.00225506975800025,
zmax=3,
ztick={\empty},
zmajorgrids,
axis x line*=bottom,
axis y line*=left,
axis z line*=left
]

\addplot3[%
surf,
shader=faceted,
draw=black,
colormap={mymap}{[1pt] rgb(0pt)=(0.2081,0.1663,0.5292); rgb(1pt)=(0.211624,0.189781,0.577676); rgb(2pt)=(0.212252,0.213771,0.626971); rgb(3pt)=(0.2081,0.2386,0.677086); rgb(4pt)=(0.195905,0.264457,0.7279); rgb(5pt)=(0.170729,0.291938,0.779248); rgb(6pt)=(0.125271,0.324243,0.830271); rgb(7pt)=(0.0591333,0.359833,0.868333); rgb(8pt)=(0.0116952,0.38751,0.881957); rgb(9pt)=(0.00595714,0.408614,0.882843); rgb(10pt)=(0.0165143,0.4266,0.878633); rgb(11pt)=(0.0328524,0.443043,0.871957); rgb(12pt)=(0.0498143,0.458571,0.864057); rgb(13pt)=(0.0629333,0.47369,0.855438); rgb(14pt)=(0.0722667,0.488667,0.8467); rgb(15pt)=(0.0779429,0.503986,0.838371); rgb(16pt)=(0.0793476,0.520024,0.831181); rgb(17pt)=(0.0749429,0.537543,0.826271); rgb(18pt)=(0.0640571,0.556986,0.823957); rgb(19pt)=(0.0487714,0.577224,0.822829); rgb(20pt)=(0.0343429,0.596581,0.819852); rgb(21pt)=(0.0265,0.6137,0.8135); rgb(22pt)=(0.0238905,0.628662,0.803762); rgb(23pt)=(0.0230905,0.641786,0.791267); rgb(24pt)=(0.0227714,0.653486,0.776757); rgb(25pt)=(0.0266619,0.664195,0.760719); rgb(26pt)=(0.0383714,0.674271,0.743552); rgb(27pt)=(0.0589714,0.683757,0.725386); rgb(28pt)=(0.0843,0.692833,0.706167); rgb(29pt)=(0.113295,0.7015,0.685857); rgb(30pt)=(0.145271,0.709757,0.664629); rgb(31pt)=(0.180133,0.717657,0.642433); rgb(32pt)=(0.217829,0.725043,0.619262); rgb(33pt)=(0.258643,0.731714,0.595429); rgb(34pt)=(0.302171,0.737605,0.571186); rgb(35pt)=(0.348167,0.742433,0.547267); rgb(36pt)=(0.395257,0.7459,0.524443); rgb(37pt)=(0.44201,0.748081,0.503314); rgb(38pt)=(0.487124,0.749062,0.483976); rgb(39pt)=(0.530029,0.749114,0.466114); rgb(40pt)=(0.570857,0.748519,0.44939); rgb(41pt)=(0.609852,0.747314,0.433686); rgb(42pt)=(0.6473,0.7456,0.4188); rgb(43pt)=(0.683419,0.743476,0.404433); rgb(44pt)=(0.71841,0.741133,0.390476); rgb(45pt)=(0.752486,0.7384,0.376814); rgb(46pt)=(0.785843,0.735567,0.363271); rgb(47pt)=(0.818505,0.732733,0.34979); rgb(48pt)=(0.850657,0.7299,0.336029); rgb(49pt)=(0.882433,0.727433,0.3217); rgb(50pt)=(0.913933,0.725786,0.306276); rgb(51pt)=(0.944957,0.726114,0.288643); rgb(52pt)=(0.973895,0.731395,0.266648); rgb(53pt)=(0.993771,0.745457,0.240348); rgb(54pt)=(0.999043,0.765314,0.216414); rgb(55pt)=(0.995533,0.786057,0.196652); rgb(56pt)=(0.988,0.8066,0.179367); rgb(57pt)=(0.978857,0.827143,0.163314); rgb(58pt)=(0.9697,0.848138,0.147452); rgb(59pt)=(0.962586,0.870514,0.1309); rgb(60pt)=(0.958871,0.8949,0.113243); rgb(61pt)=(0.959824,0.921833,0.0948381); rgb(62pt)=(0.9661,0.951443,0.0755333); rgb(63pt)=(0.9763,0.9831,0.0538)},
mesh/rows=30]
table[row sep=crcr,header=false] {0	0	0.118352773742808\\
0	0.0344827586206897	0.051402733664511\\
0	0.0689655172413793	0.0144955025456188\\
0	0.103448275862069	0.00799732940231117\\
0	0.137931034482759	0.0504457738476799\\
0	0.172413793103448	0.116305609098239\\
0	0.206896551724138	0.153409403343879\\
0	0.241379310344828	0.148710065697585\\
0	0.275862068965517	0.128988519473172\\
0	0.310344827586207	0.11814035138516\\
0	0.344827586206897	0.117198106800518\\
0	0.379310344827586	0.117538221760774\\
0	0.413793103448276	0.113677736703432\\
0	0.448275862068965	0.106044979487242\\
0	0.482758620689655	0.0965034881059803\\
0	0.517241379310345	0.0842361113476219\\
0	0.551724137931034	0.0645078897232184\\
0	0.586206896551724	0.0540427263920414\\
0	0.620689655172414	0.0744611779370698\\
0	0.655172413793103	0.126309556134795\\
0	0.689655172413793	0.168526183532657\\
0	0.724137931034483	0.16386110191874\\
0	0.758620689655172	0.125380066955824\\
0	0.793103448275862	0.0881364029872505\\
0	0.827586206896552	0.072677828640544\\
0	0.862068965517241	0.0668806979302004\\
0	0.896551724137931	0.0735481583751604\\
0	0.931034482758621	0.0944243086902338\\
0	0.96551724137931	0.128353456496835\\
0	1	0.175373514681729\\
0.0344827586206897	0	0.133051801904348\\
0.0344827586206897	0.0344827586206897	0.0669327133402057\\
0.0344827586206897	0.0689655172413793	0.0299940420564649\\
0.0344827586206897	0.103448275862069	0.0226197381365943\\
0.0344827586206897	0.137931034482759	0.0649266931826866\\
0.0344827586206897	0.172413793103448	0.129171179621575\\
0.0344827586206897	0.206896551724138	0.15827702502747\\
0.0344827586206897	0.241379310344828	0.12996101475779\\
0.0344827586206897	0.275862068965517	0.0871250092903166\\
0.0344827586206897	0.310344827586207	0.078702434972271\\
0.0344827586206897	0.344827586206897	0.114766846146786\\
0.0344827586206897	0.379310344827586	0.161304075995182\\
0.0344827586206897	0.413793103448276	0.185590452852345\\
0.0344827586206897	0.448275862068965	0.191750945245977\\
0.0344827586206897	0.482758620689655	0.186881925234792\\
0.0344827586206897	0.517241379310345	0.169995891204855\\
0.0344827586206897	0.551724137931034	0.132825062589305\\
0.0344827586206897	0.586206896551724	0.103969816338456\\
0.0344827586206897	0.620689655172414	0.118344416082904\\
0.0344827586206897	0.655172413793103	0.174665736004839\\
0.0344827586206897	0.689655172413793	0.216615116701535\\
0.0344827586206897	0.724137931034483	0.19027258156996\\
0.0344827586206897	0.758620689655172	0.112683455229781\\
0.0344827586206897	0.793103448275862	0.048776491137308\\
0.0344827586206897	0.827586206896552	0.0482088331629678\\
0.0344827586206897	0.862068965517241	0.0844715020053526\\
0.0344827586206897	0.896551724137931	0.11440568612681\\
0.0344827586206897	0.931034482758621	0.11701169982441\\
0.0344827586206897	0.96551724137931	0.107227310550705\\
0.0344827586206897	1	0.0852026183574342\\
0.0689655172413793	0	0.141486259314942\\
0.0689655172413793	0.0344827586206897	0.0776366884727265\\
0.0689655172413793	0.0689655172413793	0.0411418631588219\\
0.0689655172413793	0.103448275862069	0.0323722564806591\\
0.0689655172413793	0.137931034482759	0.0709744704071807\\
0.0689655172413793	0.172413793103448	0.129842372396177\\
0.0689655172413793	0.206896551724138	0.153020910253856\\
0.0689655172413793	0.241379310344828	0.114888844405272\\
0.0689655172413793	0.275862068965517	0.0639232654085365\\
0.0689655172413793	0.310344827586207	0.0588951673358285\\
0.0689655172413793	0.344827586206897	0.114304808497332\\
0.0689655172413793	0.379310344827586	0.18436111222759\\
0.0689655172413793	0.413793103448276	0.223262493395279\\
0.0689655172413793	0.448275862068965	0.236904595522593\\
0.0689655172413793	0.482758620689655	0.234902405485745\\
0.0689655172413793	0.517241379310345	0.216242494816764\\
0.0689655172413793	0.551724137931034	0.171067397348964\\
0.0689655172413793	0.586206896551724	0.133260409146075\\
0.0689655172413793	0.620689655172414	0.143532933042141\\
0.0689655172413793	0.655172413793103	0.199658337989242\\
0.0689655172413793	0.689655172413793	0.239334413613678\\
0.0689655172413793	0.724137931034483	0.20170280206649\\
0.0689655172413793	0.758620689655172	0.105252723452515\\
0.0689655172413793	0.793103448275862	0.0289588217363666\\
0.0689655172413793	0.827586206896552	0.0368264919521401\\
0.0689655172413793	0.862068965517241	0.0953515536560926\\
0.0689655172413793	0.896551724137931	0.137009845160383\\
0.0689655172413793	0.931034482758621	0.128727331132009\\
0.0689655172413793	0.96551724137931	0.0939815259081873\\
0.0689655172413793	1	0.032979695607701\\
0.103448275862069	0	0.143597838796122\\
0.103448275862069	0.0344827586206897	0.0834467888783432\\
0.103448275862069	0.0689655172413793	0.0478580079475322\\
0.103448275862069	0.103448275862069	0.0371570633400334\\
0.103448275862069	0.137931034482759	0.0684530485320537\\
0.103448275862069	0.172413793103448	0.118147996227669\\
0.103448275862069	0.206896551724138	0.137489871683163\\
0.103448275862069	0.241379310344828	0.103471953548799\\
0.103448275862069	0.275862068965517	0.0595039009241212\\
0.103448275862069	0.310344827586207	0.0588527872721051\\
0.103448275862069	0.344827586206897	0.115787348529833\\
0.103448275862069	0.379310344827586	0.18648189324557\\
0.103448275862069	0.413793103448276	0.226355627255373\\
0.103448275862069	0.448275862068965	0.241137143937756\\
0.103448275862069	0.482758620689655	0.240202449717333\\
0.103448275862069	0.517241379310345	0.222636246098268\\
0.103448275862069	0.551724137931034	0.178952292860701\\
0.103448275862069	0.586206896551724	0.141684248807479\\
0.103448275862069	0.620689655172414	0.149792004190184\\
0.103448275862069	0.655172413793103	0.200991870386644\\
0.103448275862069	0.689655172413793	0.236362281291979\\
0.103448275862069	0.724137931034483	0.197938578164012\\
0.103448275862069	0.758620689655172	0.103093616389193\\
0.103448275862069	0.793103448275862	0.0288435117516358\\
0.103448275862069	0.827586206896552	0.0386179472720639\\
0.103448275862069	0.862068965517241	0.0993892475661005\\
0.103448275862069	0.896551724137931	0.141100612127285\\
0.103448275862069	0.931034482758621	0.129389752354533\\
0.103448275862069	0.96551724137931	0.0886376721766004\\
0.103448275862069	1	0.0190532064299762\\
0.137931034482759	0	0.135299639717406\\
0.137931034482759	0.0344827586206897	0.0795022777920944\\
0.137931034482759	0.0689655172413793	0.044283279563702\\
0.137931034482759	0.103448275862069	0.0298745753981053\\
0.137931034482759	0.137931034482759	0.0475514112902892\\
0.137931034482759	0.172413793103448	0.0818047108644439\\
0.137931034482759	0.206896551724138	0.100802849782935\\
0.137931034482759	0.241379310344828	0.0939216749227482\\
0.137931034482759	0.275862068965517	0.0820647164882248\\
0.137931034482759	0.310344827586207	0.0877531598607738\\
0.137931034482759	0.344827586206897	0.11728571045659\\
0.137931034482759	0.379310344827586	0.151568743274301\\
0.137931034482759	0.413793103448276	0.171083872234316\\
0.137931034482759	0.448275862068965	0.178658786914007\\
0.137931034482759	0.482758620689655	0.17754707317075\\
0.137931034482759	0.517241379310345	0.165525211592339\\
0.137931034482759	0.551724137931034	0.136682274411662\\
0.137931034482759	0.586206896551724	0.11294257356728\\
0.137931034482759	0.620689655172414	0.12040120321285\\
0.137931034482759	0.655172413793103	0.157604105553331\\
0.137931034482759	0.689655172413793	0.184755257128543\\
0.137931034482759	0.724137931034483	0.16369761558911\\
0.137931034482759	0.758620689655172	0.106398430074816\\
0.137931034482759	0.793103448275862	0.059533717062264\\
0.137931034482759	0.827586206896552	0.059517665505532\\
0.137931034482759	0.862068965517241	0.0870452414074104\\
0.137931034482759	0.896551724137931	0.108051058217735\\
0.137931034482759	0.931034482758621	0.105915246617677\\
0.137931034482759	0.96551724137931	0.0924492499579433\\
0.137931034482759	1	0.0677654529819487\\
0.172413793103448	0	0.122285010826383\\
0.172413793103448	0.0344827586206897	0.0725778595994218\\
0.172413793103448	0.0689655172413793	0.0385859336747086\\
0.172413793103448	0.103448275862069	0.0204228897614719\\
0.172413793103448	0.137931034482759	0.0219459050668695\\
0.172413793103448	0.172413793103448	0.0379333147093765\\
0.172413793103448	0.206896551724138	0.0581271638401765\\
0.172413793103448	0.241379310344828	0.0887386657290751\\
0.172413793103448	0.275862068965517	0.120193907782249\\
0.172413793103448	0.310344827586207	0.132818366787452\\
0.172413793103448	0.344827586206897	0.121493303632934\\
0.172413793103448	0.379310344827586	0.102051892496468\\
0.172413793103448	0.413793103448276	0.0905852818114294\\
0.172413793103448	0.448275862068965	0.0853945618560068\\
0.172413793103448	0.482758620689655	0.0820847498025457\\
0.172413793103448	0.517241379310345	0.0778530505902464\\
0.172413793103448	0.551724137931034	0.0718361994951512\\
0.172413793103448	0.586206896551724	0.06974580612522\\
0.172413793103448	0.620689655172414	0.0786618708391407\\
0.172413793103448	0.655172413793103	0.0988474776934657\\
0.172413793103448	0.689655172413793	0.116487682454615\\
0.172413793103448	0.724137931034483	0.120280163677554\\
0.172413793103448	0.758620689655172	0.114906102016436\\
0.172413793103448	0.793103448275862	0.105565901196989\\
0.172413793103448	0.827586206896552	0.0912641690783366\\
0.172413793103448	0.862068965517241	0.0716186414684238\\
0.172413793103448	0.896551724137931	0.0638226755044452\\
0.172413793103448	0.931034482758621	0.0765420986199597\\
0.172413793103448	0.96551724137931	0.103678432988704\\
0.172413793103448	1	0.145208765661568\\
0.206896551724138	0	0.116090220303241\\
0.206896551724138	0.0344827586206897	0.0766127372223685\\
0.206896551724138	0.0689655172413793	0.0476740654333492\\
0.206896551724138	0.103448275862069	0.0292931682516207\\
0.206896551724138	0.137931034482759	0.0197683366499497\\
0.206896551724138	0.172413793103448	0.0215721727258632\\
0.206896551724138	0.206896551724138	0.0404833361177542\\
0.206896551724138	0.241379310344828	0.0932245451357916\\
0.206896551724138	0.275862068965517	0.150961116888334\\
0.206896551724138	0.310344827586207	0.168359129949028\\
0.206896551724138	0.344827586206897	0.134183100038607\\
0.206896551724138	0.379310344827586	0.083792824115459\\
0.206896551724138	0.413793103448276	0.0523261430527321\\
0.206896551724138	0.448275862068965	0.0340996162450368\\
0.206896551724138	0.482758620689655	0.0246663428897703\\
0.206896551724138	0.517241379310345	0.026022478695585\\
0.206896551724138	0.551724137931034	0.0402985906051787\\
0.206896551724138	0.586206896551724	0.0585398194279353\\
0.206896551724138	0.620689655172414	0.0725327219756704\\
0.206896551724138	0.655172413793103	0.0852467789191759\\
0.206896551724138	0.689655172413793	0.0975534052752696\\
0.206896551724138	0.724137931034483	0.111861574828819\\
0.206896551724138	0.758620689655172	0.128605755173137\\
0.206896551724138	0.793103448275862	0.135731869365111\\
0.206896551724138	0.827586206896552	0.11730617920266\\
0.206896551724138	0.862068965517241	0.0805839615020531\\
0.206896551724138	0.896551724137931	0.0617445306939657\\
0.206896551724138	0.931034482758621	0.0788016720692962\\
0.206896551724138	0.96551724137931	0.119030358627723\\
0.206896551724138	1	0.182354472647724\\
0.241379310344828	0	0.119133202902512\\
0.241379310344828	0.0344827586206897	0.0986798373625899\\
0.241379310344828	0.0689655172413793	0.0820910029051492\\
0.241379310344828	0.103448275862069	0.0693127740613557\\
0.241379310344828	0.137931034482759	0.0551617451716015\\
0.241379310344828	0.172413793103448	0.046896474248146\\
0.241379310344828	0.206896551724138	0.0600451064476429\\
0.241379310344828	0.241379310344828	0.113134061710056\\
0.241379310344828	0.275862068965517	0.173392785756001\\
0.241379310344828	0.310344827586207	0.19299611329098\\
0.241379310344828	0.344827586206897	0.162759487531463\\
0.241379310344828	0.379310344827586	0.114454486625598\\
0.241379310344828	0.413793103448276	0.0766780206645414\\
0.241379310344828	0.448275862068965	0.0392394632928352\\
0.241379310344828	0.482758620689655	0.0128885797746676\\
0.241379310344828	0.517241379310345	0.0170798070895314\\
0.241379310344828	0.551724137931034	0.0538113113324839\\
0.241379310344828	0.586206896551724	0.0974309010010766\\
0.241379310344828	0.620689655172414	0.126606669659806\\
0.241379310344828	0.655172413793103	0.149982413679464\\
0.241379310344828	0.689655172413793	0.165330538074527\\
0.241379310344828	0.724137931034483	0.167543058946031\\
0.241379310344828	0.758620689655172	0.157638058082856\\
0.241379310344828	0.793103448275862	0.145969749616662\\
0.241379310344828	0.827586206896552	0.137902209808517\\
0.241379310344828	0.862068965517241	0.130034216615655\\
0.241379310344828	0.896551724137931	0.127179784890829\\
0.241379310344828	0.931034482758621	0.131853196541028\\
0.241379310344828	0.96551724137931	0.142296286139616\\
0.241379310344828	1	0.158509066913343\\
0.275862068965517	0	0.126228540519939\\
0.275862068965517	0.0344827586206897	0.125438909587913\\
0.275862068965517	0.0689655172413793	0.122308570212351\\
0.275862068965517	0.103448275862069	0.11673109219214\\
0.275862068965517	0.137931034482759	0.101429535441309\\
0.275862068965517	0.172413793103448	0.0865190297048114\\
0.275862068965517	0.206896551724138	0.0932079766143626\\
0.275862068965517	0.241379310344828	0.138175964016861\\
0.275862068965517	0.275862068965517	0.191269372294377\\
0.275862068965517	0.310344827586207	0.211454868457161\\
0.275862068965517	0.344827586206897	0.194072698454456\\
0.275862068965517	0.379310344827586	0.159702271792519\\
0.275862068965517	0.413793103448276	0.122695742896972\\
0.275862068965517	0.448275862068965	0.0694473349873049\\
0.275862068965517	0.482758620689655	0.0272678991703979\\
0.275862068965517	0.517241379310345	0.0328902295487476\\
0.275862068965517	0.551724137931034	0.0872071454220341\\
0.275862068965517	0.586206896551724	0.151275318017284\\
0.275862068965517	0.620689655172414	0.194544523181365\\
0.275862068965517	0.655172413793103	0.230919006137001\\
0.275862068965517	0.689655172413793	0.250024298704774\\
0.275862068965517	0.724137931034483	0.233649813085982\\
0.275862068965517	0.758620689655172	0.184720933595425\\
0.275862068965517	0.793103448275862	0.146082764145614\\
0.275862068965517	0.827586206896552	0.154141215851454\\
0.275862068965517	0.862068965517241	0.190011014739296\\
0.275862068965517	0.896551724137931	0.21198657151101\\
0.275862068965517	0.931034482758621	0.199586131600206\\
0.275862068965517	0.96551724137931	0.167323913616879\\
0.275862068965517	1	0.11531363527303\\
0.310344827586207	0	0.134191593003404\\
0.310344827586207	0.0344827586206897	0.140947852481807\\
0.310344827586207	0.0689655172413793	0.143200423201141\\
0.310344827586207	0.103448275862069	0.140829155025793\\
0.310344827586207	0.137931034482759	0.126170735049256\\
0.310344827586207	0.172413793103448	0.109857636323338\\
0.310344827586207	0.206896551724138	0.114071736753269\\
0.310344827586207	0.241379310344828	0.153981720709773\\
0.310344827586207	0.275862068965517	0.201904679934877\\
0.310344827586207	0.310344827586207	0.221231656600534\\
0.310344827586207	0.344827586206897	0.208290518105646\\
0.310344827586207	0.379310344827586	0.180441613318679\\
0.310344827586207	0.413793103448276	0.149101861838007\\
0.310344827586207	0.448275862068965	0.102008290204468\\
0.310344827586207	0.482758620689655	0.0646014178841477\\
0.310344827586207	0.517241379310345	0.0708287350915032\\
0.310344827586207	0.551724137931034	0.121731985639464\\
0.310344827586207	0.586206896551724	0.180578734195974\\
0.310344827586207	0.620689655172414	0.218402481573495\\
0.310344827586207	0.655172413793103	0.24826275654732\\
0.310344827586207	0.689655172413793	0.260866119913882\\
0.310344827586207	0.724137931034483	0.237596134959999\\
0.310344827586207	0.758620689655172	0.180626695129957\\
0.310344827586207	0.793103448275862	0.139991816825118\\
0.310344827586207	0.827586206896552	0.16184930510632\\
0.310344827586207	0.862068965517241	0.222894303643116\\
0.310344827586207	0.896551724137931	0.258916020898569\\
0.310344827586207	0.931034482758621	0.238224839796621\\
0.310344827586207	0.96551724137931	0.183252177243882\\
0.310344827586207	1	0.0941590613186511\\
0.344827586206897	0	0.150477856534634\\
0.344827586206897	0.0344827586206897	0.147105934323232\\
0.344827586206897	0.0689655172413793	0.143128521788821\\
0.344827586206897	0.103448275862069	0.138477107669046\\
0.344827586206897	0.137931034482759	0.128200366614058\\
0.344827586206897	0.172413793103448	0.119190307335015\\
0.344827586206897	0.206896551724138	0.125964981064221\\
0.344827586206897	0.241379310344828	0.161303291141811\\
0.344827586206897	0.275862068965517	0.202256632464377\\
0.344827586206897	0.310344827586207	0.216560907207304\\
0.344827586206897	0.344827586206897	0.196873394602085\\
0.344827586206897	0.379310344827586	0.167026363817963\\
0.344827586206897	0.413793103448276	0.150327175156199\\
0.344827586206897	0.448275862068965	0.143233653946665\\
0.344827586206897	0.482758620689655	0.142214923095711\\
0.344827586206897	0.517241379310345	0.148648342706845\\
0.344827586206897	0.551724137931034	0.165687728404757\\
0.344827586206897	0.586206896551724	0.181478841741684\\
0.344827586206897	0.620689655172414	0.184247983088685\\
0.344827586206897	0.655172413793103	0.17715716416026\\
0.344827586206897	0.689655172413793	0.166395367467083\\
0.344827586206897	0.724137931034483	0.151873670264035\\
0.344827586206897	0.758620689655172	0.130722407849885\\
0.344827586206897	0.793103448275862	0.124913683621353\\
0.344827586206897	0.827586206896552	0.162040678672934\\
0.344827586206897	0.862068965517241	0.229328437068657\\
0.344827586206897	0.896551724137931	0.267832882506616\\
0.344827586206897	0.931034482758621	0.248230633267016\\
0.344827586206897	0.96551724137931	0.191238030579709\\
0.344827586206897	1	0.0969802734102918\\
0.379310344827586	0	0.16042688399714\\
0.379310344827586	0.0344827586206897	0.144310062511479\\
0.379310344827586	0.0689655172413793	0.13236664470501\\
0.379310344827586	0.103448275862069	0.124584781215994\\
0.379310344827586	0.137931034482759	0.11883992083167\\
0.379310344827586	0.172413793103448	0.118132227418757\\
0.379310344827586	0.206896551724138	0.129035959826738\\
0.379310344827586	0.241379310344828	0.162499715257182\\
0.379310344827586	0.275862068965517	0.199388723416372\\
0.379310344827586	0.310344827586207	0.210128721068031\\
0.379310344827586	0.344827586206897	0.183027230292638\\
0.379310344827586	0.379310344827586	0.150239881644518\\
0.379310344827586	0.413793103448276	0.149266338283698\\
0.379310344827586	0.448275862068965	0.18620592944677\\
0.379310344827586	0.482758620689655	0.224920834173691\\
0.379310344827586	0.517241379310345	0.230111398610574\\
0.379310344827586	0.551724137931034	0.206920530499779\\
0.379310344827586	0.586206896551724	0.172816169892325\\
0.379310344827586	0.620689655172414	0.136934296189796\\
0.379310344827586	0.655172413793103	0.0911799170470698\\
0.379310344827586	0.689655172413793	0.0569447630265378\\
0.379310344827586	0.724137931034483	0.0536112184858376\\
0.379310344827586	0.758620689655172	0.0734655011473275\\
0.379310344827586	0.793103448275862	0.105562286154651\\
0.379310344827586	0.827586206896552	0.153732636498086\\
0.379310344827586	0.862068965517241	0.218244564891265\\
0.379310344827586	0.896551724137931	0.254658876828661\\
0.379310344827586	0.931034482758621	0.240601014512998\\
0.379310344827586	0.96551724137931	0.191834371121725\\
0.379310344827586	1	0.108428529308722\\
0.413793103448276	0	0.140773770944012\\
0.413793103448276	0.0344827586206897	0.127668060155434\\
0.413793103448276	0.0689655172413793	0.117706813880564\\
0.413793103448276	0.103448275862069	0.110852883204933\\
0.413793103448276	0.137931034482759	0.103391931249166\\
0.413793103448276	0.172413793103448	0.100528955946611\\
0.413793103448276	0.206896551724138	0.113401080415469\\
0.413793103448276	0.241379310344828	0.155366425103497\\
0.413793103448276	0.275862068965517	0.202805407823817\\
0.413793103448276	0.310344827586207	0.22048958124131\\
0.413793103448276	0.344827586206897	0.195485727940309\\
0.413793103448276	0.379310344827586	0.164044549334073\\
0.413793103448276	0.413793103448276	0.168369080235477\\
0.413793103448276	0.448275862068965	0.217133450679754\\
0.413793103448276	0.482758620689655	0.264967164393782\\
0.413793103448276	0.517241379310345	0.265846360805662\\
0.413793103448276	0.551724137931034	0.224411134496176\\
0.413793103448276	0.586206896551724	0.170326623146438\\
0.413793103448276	0.620689655172414	0.123223070403615\\
0.413793103448276	0.655172413793103	0.0714884310356493\\
0.413793103448276	0.689655172413793	0.0341995311834525\\
0.413793103448276	0.724137931034483	0.0304634015637129\\
0.413793103448276	0.758620689655172	0.0532998503806878\\
0.413793103448276	0.793103448275862	0.0873272196526517\\
0.413793103448276	0.827586206896552	0.131460871912399\\
0.413793103448276	0.862068965517241	0.188337546065579\\
0.413793103448276	0.896551724137931	0.222521982491845\\
0.413793103448276	0.931034482758621	0.216110401136467\\
0.413793103448276	0.96551724137931	0.181714912045506\\
0.413793103448276	1	0.119390296128844\\
0.448275862068965	0	0.0800874912024652\\
0.448275862068965	0.0344827586206897	0.0912501535985415\\
0.448275862068965	0.0689655172413793	0.0961355464772372\\
0.448275862068965	0.103448275862069	0.0945554574995199\\
0.448275862068965	0.137931034482759	0.0742759962955863\\
0.448275862068965	0.172413793103448	0.0522798688238859\\
0.448275862068965	0.206896551724138	0.0639374490898043\\
0.448275862068965	0.241379310344828	0.132378353936572\\
0.448275862068965	0.275862068965517	0.215273058098988\\
0.448275862068965	0.310344827586207	0.256500351000821\\
0.448275862068965	0.344827586206897	0.244738874287618\\
0.448275862068965	0.379310344827586	0.218485418589269\\
0.448275862068965	0.413793103448276	0.216436398495497\\
0.448275862068965	0.448275862068965	0.241948438286783\\
0.448275862068965	0.482758620689655	0.265212240048605\\
0.448275862068965	0.517241379310345	0.257199614323913\\
0.448275862068965	0.551724137931034	0.219050062082742\\
0.448275862068965	0.586206896551724	0.175461540849934\\
0.448275862068965	0.620689655172414	0.146973457952632\\
0.448275862068965	0.655172413793103	0.126836359658158\\
0.448275862068965	0.689655172413793	0.109825013462124\\
0.448275862068965	0.724137931034483	0.0905540814116926\\
0.448275862068965	0.758620689655172	0.0697495257583777\\
0.448275862068965	0.793103448275862	0.0615561984553503\\
0.448275862068965	0.827586206896552	0.0818532040764363\\
0.448275862068965	0.862068965517241	0.123127956469036\\
0.448275862068965	0.896551724137931	0.154889825917981\\
0.448275862068965	0.931034482758621	0.162005198815952\\
0.448275862068965	0.96551724137931	0.15520367076546\\
0.448275862068965	1	0.134572410803068\\
0.482758620689655	0	0.0208535523785628\\
0.482758620689655	0.0344827586206897	0.0531348507316472\\
0.482758620689655	0.0689655172413793	0.0717217434935502\\
0.482758620689655	0.103448275862069	0.0763108217811279\\
0.482758620689655	0.137931034482759	0.0482409251857022\\
0.482758620689655	0.172413793103448	0.0133770107651516\\
0.482758620689655	0.206896551724138	0.0253477517582603\\
0.482758620689655	0.241379310344828	0.114282345347386\\
0.482758620689655	0.275862068965517	0.22436867061512\\
0.482758620689655	0.310344827586207	0.284940210099013\\
0.482758620689655	0.344827586206897	0.287719499727242\\
0.482758620689655	0.379310344827586	0.268770848509122\\
0.482758620689655	0.413793103448276	0.258246943473236\\
0.482758620689655	0.448275862068965	0.252837229292159\\
0.482758620689655	0.482758620689655	0.244525910785513\\
0.482758620689655	0.517241379310345	0.227856825074865\\
0.482758620689655	0.551724137931034	0.200088919714735\\
0.482758620689655	0.586206896551724	0.175771697622884\\
0.482758620689655	0.620689655172414	0.172397896839844\\
0.482758620689655	0.655172413793103	0.189787600529866\\
0.482758620689655	0.689655172413793	0.197523102006273\\
0.482758620689655	0.724137931034483	0.164511148250767\\
0.482758620689655	0.758620689655172	0.0995733403557704\\
0.482758620689655	0.793103448275862	0.0477472240421912\\
0.482758620689655	0.827586206896552	0.0426330702589863\\
0.482758620689655	0.862068965517241	0.0661033873964723\\
0.482758620689655	0.896551724137931	0.0929387696444077\\
0.482758620689655	0.931034482758621	0.110950960535838\\
0.482758620689655	0.96551724137931	0.128859924241815\\
0.482758620689655	1	0.146782483385778\\
0.517241379310345	0	0.0111431191652372\\
0.517241379310345	0.0344827586206897	0.0322507373121429\\
0.517241379310345	0.0689655172413793	0.0471286796685421\\
0.517241379310345	0.103448275862069	0.0555962343794589\\
0.517241379310345	0.137931034482759	0.0459675391377862\\
0.517241379310345	0.172413793103448	0.0344622631716348\\
0.517241379310345	0.206896551724138	0.054826185813523\\
0.517241379310345	0.241379310344828	0.128431866394626\\
0.517241379310345	0.275862068965517	0.216074823496804\\
0.517241379310345	0.310344827586207	0.266877139097283\\
0.517241379310345	0.344827586206897	0.276706388381732\\
0.517241379310345	0.379310344827586	0.267706415258231\\
0.517241379310345	0.413793103448276	0.255742190513908\\
0.517241379310345	0.448275862068965	0.235533546202566\\
0.517241379310345	0.482758620689655	0.211925853372648\\
0.517241379310345	0.517241379310345	0.191881017151356\\
0.517241379310345	0.551724137931034	0.170899272743131\\
0.517241379310345	0.586206896551724	0.157217668479808\\
0.517241379310345	0.620689655172414	0.164383644078694\\
0.517241379310345	0.655172413793103	0.194533707530261\\
0.517241379310345	0.689655172413793	0.213538286628372\\
0.517241379310345	0.724137931034483	0.187388090581582\\
0.517241379310345	0.758620689655172	0.126453665894718\\
0.517241379310345	0.793103448275862	0.0744757139635913\\
0.517241379310345	0.827586206896552	0.062153831546707\\
0.517241379310345	0.862068965517241	0.0725572877651025\\
0.517241379310345	0.896551724137931	0.089032145944455\\
0.517241379310345	0.931034482758621	0.103664936192317\\
0.517241379310345	0.96551724137931	0.122144980662457\\
0.517241379310345	1	0.144564564552818\\
0.551724137931034	0	0.0513455014760907\\
0.551724137931034	0.0344827586206897	0.0232287273136016\\
0.551724137931034	0.0689655172413793	0.0142468731587392\\
0.551724137931034	0.103448275862069	0.0246237478529873\\
0.551724137931034	0.137931034482759	0.065444535595246\\
0.551724137931034	0.172413793103448	0.121451336919288\\
0.551724137931034	0.206896551724138	0.161565119076051\\
0.551724137931034	0.241379310344828	0.179916703873285\\
0.551724137931034	0.275862068965517	0.189231575415788\\
0.551724137931034	0.310344827586207	0.198825830532454\\
0.551724137931034	0.344827586206897	0.211413860601807\\
0.551724137931034	0.379310344827586	0.218775710723893\\
0.551724137931034	0.413793103448276	0.211307753053875\\
0.551724137931034	0.448275862068965	0.185309517802464\\
0.551724137931034	0.482758620689655	0.155475734172616\\
0.551724137931034	0.517241379310345	0.136045474625466\\
0.551724137931034	0.551724137931034	0.122261657956169\\
0.551724137931034	0.586206896551724	0.116052236091912\\
0.551724137931034	0.620689655172414	0.1233405754272\\
0.551724137931034	0.655172413793103	0.145565532161267\\
0.551724137931034	0.689655172413793	0.165563687048783\\
0.551724137931034	0.724137931034483	0.168257266506241\\
0.551724137931034	0.758620689655172	0.159370028664406\\
0.551724137931034	0.793103448275862	0.149863062847754\\
0.551724137931034	0.827586206896552	0.147168461748072\\
0.551724137931034	0.862068965517241	0.147171040697491\\
0.551724137931034	0.896551724137931	0.145792801922829\\
0.551724137931034	0.931034482758621	0.141088957842194\\
0.551724137931034	0.96551724137931	0.134437210891454\\
0.551724137931034	1	0.125848090618787\\
0.586206896551724	0	0.0921531204250033\\
0.586206896551724	0.0344827586206897	0.0252624179150256\\
0.586206896551724	0.0689655172413793	-0.00225506975800025\\
0.586206896551724	0.103448275862069	0.0101611743011517\\
0.586206896551724	0.137931034482759	0.0928079710377732\\
0.586206896551724	0.172413793103448	0.203859018350262\\
0.586206896551724	0.206896551724138	0.257496044172491\\
0.586206896551724	0.241379310344828	0.224279632033306\\
0.586206896551724	0.275862068965517	0.161737893855388\\
0.586206896551724	0.310344827586207	0.131391111415659\\
0.586206896551724	0.344827586206897	0.140440822304741\\
0.586206896551724	0.379310344827586	0.15735147163175\\
0.586206896551724	0.413793103448276	0.154890356984089\\
0.586206896551724	0.448275862068965	0.132315247016041\\
0.586206896551724	0.482758620689655	0.10635821541112\\
0.586206896551724	0.517241379310345	0.0908450432405399\\
0.586206896551724	0.551724137931034	0.0818551401436376\\
0.586206896551724	0.586206896551724	0.0789137629613721\\
0.586206896551724	0.620689655172414	0.0829506482244801\\
0.586206896551724	0.655172413793103	0.0934713825639663\\
0.586206896551724	0.689655172413793	0.11132323477095\\
0.586206896551724	0.724137931034483	0.141417196028876\\
0.586206896551724	0.758620689655172	0.1843780895814\\
0.586206896551724	0.793103448275862	0.218029972458396\\
0.586206896551724	0.827586206896552	0.226659372090935\\
0.586206896551724	0.862068965517241	0.218975919970999\\
0.586206896551724	0.896551724137931	0.202111011659186\\
0.586206896551724	0.931034482758621	0.179401602479927\\
0.586206896551724	0.96551724137931	0.148401393859139\\
0.586206896551724	1	0.109043462591841\\
0.620689655172414	0	0.083629490035255\\
0.620689655172414	0.0344827586206897	0.0422572105385316\\
0.620689655172414	0.0689655172413793	0.0302687934456584\\
0.620689655172414	0.103448275862069	0.0481219864791208\\
0.620689655172414	0.137931034482759	0.121379689685225\\
0.620689655172414	0.172413793103448	0.214713904373976\\
0.620689655172414	0.206896551724138	0.255586944774854\\
0.620689655172414	0.241379310344828	0.217772845790577\\
0.620689655172414	0.275862068965517	0.152139308984681\\
0.620689655172414	0.310344827586207	0.113891749046325\\
0.620689655172414	0.344827586206897	0.105740569612684\\
0.620689655172414	0.379310344827586	0.107180544903217\\
0.620689655172414	0.413793103448276	0.105104254369062\\
0.620689655172414	0.448275862068965	0.102129254723396\\
0.620689655172414	0.482758620689655	0.0986823827623906\\
0.620689655172414	0.517241379310345	0.0921748312129363\\
0.620689655172414	0.551724137931034	0.0811019591059577\\
0.620689655172414	0.586206896551724	0.0730749214532372\\
0.620689655172414	0.620689655172414	0.0744917871586565\\
0.620689655172414	0.655172413793103	0.0824042653600716\\
0.620689655172414	0.689655172413793	0.0993134817727167\\
0.620689655172414	0.724137931034483	0.13278124322143\\
0.620689655172414	0.758620689655172	0.182876530715809\\
0.620689655172414	0.793103448275862	0.222133341693245\\
0.620689655172414	0.827586206896552	0.230421937722978\\
0.620689655172414	0.862068965517241	0.218751516476104\\
0.620689655172414	0.896551724137931	0.199117597300038\\
0.620689655172414	0.931034482758621	0.177242584627933\\
0.620689655172414	0.96551724137931	0.148986914654684\\
0.620689655172414	1	0.114268734720389\\
0.655172413793103	0	0.0314253136942762\\
0.655172413793103	0.0344827586206897	0.0763908385041801\\
0.655172413793103	0.0689655172413793	0.112785908076356\\
0.655172413793103	0.103448275862069	0.140569087870519\\
0.655172413793103	0.137931034482759	0.159180323811766\\
0.655172413793103	0.172413793103448	0.169316850163282\\
0.655172413793103	0.206896551724138	0.172449832562363\\
0.655172413793103	0.241379310344828	0.16920017814856\\
0.655172413793103	0.275862068965517	0.15829332447146\\
0.655172413793103	0.310344827586207	0.13699044575553\\
0.655172413793103	0.344827586206897	0.0943832624203017\\
0.655172413793103	0.379310344827586	0.0538034964733116\\
0.655172413793103	0.413793103448276	0.0479017954845307\\
0.655172413793103	0.448275862068965	0.0839593443497137\\
0.655172413793103	0.482758620689655	0.125512627163419\\
0.655172413793103	0.517241379310345	0.135051781943033\\
0.655172413793103	0.551724137931034	0.115713703734793\\
0.655172413793103	0.586206896551724	0.0939863711862191\\
0.655172413793103	0.620689655172414	0.0915581018170841\\
0.655172413793103	0.655172413793103	0.101831401712543\\
0.655172413793103	0.689655172413793	0.11689951911774\\
0.655172413793103	0.724137931034483	0.134169373461015\\
0.655172413793103	0.758620689655172	0.156402482114773\\
0.655172413793103	0.793103448275862	0.172223344974704\\
0.655172413793103	0.827586206896552	0.171494023340044\\
0.655172413793103	0.862068965517241	0.159431095510098\\
0.655172413793103	0.896551724137931	0.147983418179457\\
0.655172413793103	0.931034482758621	0.143016782976236\\
0.655172413793103	0.96551724137931	0.140350513614423\\
0.655172413793103	1	0.139937841899004\\
0.689655172413793	0	-0.000654299197033799\\
0.689655172413793	0.0344827586206897	0.110374075306438\\
0.689655172413793	0.0689655172413793	0.181689874298995\\
0.689655172413793	0.103448275862069	0.212825819356844\\
0.689655172413793	0.137931034482759	0.180571981461278\\
0.689655172413793	0.172413793103448	0.116878257377803\\
0.689655172413793	0.206896551724138	0.0874220591076256\\
0.689655172413793	0.241379310344828	0.116302744272925\\
0.689655172413793	0.275862068965517	0.156710928024571\\
0.689655172413793	0.310344827586207	0.155356566738263\\
0.689655172413793	0.344827586206897	0.0912621436944691\\
0.689655172413793	0.379310344827586	0.0222515149462735\\
0.689655172413793	0.413793103448276	0.0156686140726228\\
0.689655172413793	0.448275862068965	0.0822880422308039\\
0.689655172413793	0.482758620689655	0.157766395313512\\
0.689655172413793	0.517241379310345	0.177979931067231\\
0.689655172413793	0.551724137931034	0.149015557951788\\
0.689655172413793	0.586206896551724	0.113701424286519\\
0.689655172413793	0.620689655172414	0.107680612162802\\
0.689655172413793	0.655172413793103	0.121388901383562\\
0.689655172413793	0.689655172413793	0.135739022479949\\
0.689655172413793	0.724137931034483	0.137149157684503\\
0.689655172413793	0.758620689655172	0.13137121224613\\
0.689655172413793	0.793103448275862	0.123877182866325\\
0.689655172413793	0.827586206896552	0.115583680845252\\
0.689655172413793	0.862068965517241	0.105441344311221\\
0.689655172413793	0.896551724137931	0.102698748369799\\
0.689655172413793	0.931034482758621	0.112041256744706\\
0.689655172413793	0.96551724137931	0.130178794361701\\
0.689655172413793	1	0.157103235093215\\
0.724137931034483	0	0.0483124128127205\\
0.724137931034483	0.0344827586206897	0.12834468210693\\
0.724137931034483	0.0689655172413793	0.176303008274532\\
0.724137931034483	0.103448275862069	0.19179154742672\\
0.724137931034483	0.137931034482759	0.154684281488855\\
0.724137931034483	0.172413793103448	0.0927087555332406\\
0.724137931034483	0.206896551724138	0.0628045598137385\\
0.724137931034483	0.241379310344828	0.084941365833947\\
0.724137931034483	0.275862068965517	0.120137586165862\\
0.724137931034483	0.310344827586207	0.125107249063875\\
0.724137931034483	0.344827586206897	0.0848186141262987\\
0.724137931034483	0.379310344827586	0.0421094545662071\\
0.724137931034483	0.413793103448276	0.0464971750295245\\
0.724137931034483	0.448275862068965	0.108617373371896\\
0.724137931034483	0.482758620689655	0.174133952080026\\
0.724137931034483	0.517241379310345	0.186616789007142\\
0.724137931034483	0.551724137931034	0.149131273277308\\
0.724137931034483	0.586206896551724	0.106993582552005\\
0.724137931034483	0.620689655172414	0.0997334040681459\\
0.724137931034483	0.655172413793103	0.117368708388521\\
0.724137931034483	0.689655172413793	0.135861715476376\\
0.724137931034483	0.724137931034483	0.137415960661708\\
0.724137931034483	0.758620689655172	0.129368657856873\\
0.724137931034483	0.793103448275862	0.120409381426136\\
0.724137931034483	0.827586206896552	0.114687450327936\\
0.724137931034483	0.862068965517241	0.10958773210144\\
0.724137931034483	0.896551724137931	0.108599511108912\\
0.724137931034483	0.931034482758621	0.113550087048051\\
0.724137931034483	0.96551724137931	0.123165071565589\\
0.724137931034483	1	0.137446482850613\\
0.758620689655172	0	0.159333758468723\\
0.758620689655172	0.0344827586206897	0.136030769379348\\
0.758620689655172	0.0689655172413793	0.116354469794134\\
0.758620689655172	0.103448275862069	0.100332728887497\\
0.758620689655172	0.137931034482759	0.0889762708874563\\
0.758620689655172	0.172413793103448	0.0809121538236621\\
0.758620689655172	0.206896551724138	0.0731855498802941\\
0.758620689655172	0.241379310344828	0.0623124176706406\\
0.758620689655172	0.275862068965517	0.054508064845241\\
0.758620689655172	0.310344827586207	0.0589769264835494\\
0.758620689655172	0.344827586206897	0.0788023646111296\\
0.758620689655172	0.379310344827586	0.105190351911959\\
0.758620689655172	0.413793103448276	0.130306311424315\\
0.758620689655172	0.448275862068965	0.161861888907938\\
0.758620689655172	0.482758620689655	0.184164051020787\\
0.758620689655172	0.517241379310345	0.174310697174401\\
0.758620689655172	0.551724137931034	0.127540310479684\\
0.758620689655172	0.586206896551724	0.0822547639118116\\
0.758620689655172	0.620689655172414	0.0755827158974568\\
0.758620689655172	0.655172413793103	0.099030609344905\\
0.758620689655172	0.689655172413793	0.126048259544999\\
0.758620689655172	0.724137931034483	0.13741422918517\\
0.758620689655172	0.758620689655172	0.141644308086615\\
0.758620689655172	0.793103448275862	0.143759928739083\\
0.758620689655172	0.827586206896552	0.147761778592933\\
0.758620689655172	0.862068965517241	0.151652074840694\\
0.758620689655172	0.896551724137931	0.149018408031038\\
0.758620689655172	0.931034482758621	0.136678848921844\\
0.758620689655172	0.96551724137931	0.116857055297064\\
0.758620689655172	1	0.0895521991027827\\
0.793103448275862	0	0.25474086410691\\
0.793103448275862	0.0344827586206897	0.147715814278392\\
0.793103448275862	0.0689655172413793	0.0730352368355962\\
0.793103448275862	0.103448275862069	0.0310602654698972\\
0.793103448275862	0.137931034482759	0.0392354837080766\\
0.793103448275862	0.172413793103448	0.0735716231130971\\
0.793103448275862	0.206896551724138	0.0846033125719258\\
0.793103448275862	0.241379310344828	0.0514549995552281\\
0.793103448275862	0.275862068965517	0.0140514918703657\\
0.793103448275862	0.310344827586207	0.0201387058979995\\
0.793103448275862	0.344827586206897	0.0868847867570768\\
0.793103448275862	0.379310344827586	0.165952897472322\\
0.793103448275862	0.413793103448276	0.204946878109132\\
0.793103448275862	0.448275862068965	0.20839871170025\\
0.793103448275862	0.482758620689655	0.192976978121204\\
0.793103448275862	0.517241379310345	0.164793049774034\\
0.793103448275862	0.551724137931034	0.113186519972822\\
0.793103448275862	0.586206896551724	0.0676586027952644\\
0.793103448275862	0.620689655172414	0.0597408944152269\\
0.793103448275862	0.655172413793103	0.0824179827053741\\
0.793103448275862	0.689655172413793	0.112359916059785\\
0.793103448275862	0.724137931034483	0.134260850313368\\
0.793103448275862	0.758620689655172	0.156028222350882\\
0.793103448275862	0.793103448275862	0.173313401344698\\
0.793103448275862	0.827586206896552	0.185224570467607\\
0.793103448275862	0.862068965517241	0.192836309894573\\
0.793103448275862	0.896551724137931	0.185397262748069\\
0.793103448275862	0.931034482758621	0.157450368683145\\
0.793103448275862	0.96551724137931	0.112781669592997\\
0.793103448275862	1	0.0513733954824769\\
0.827586206896552	0	0.282830875307032\\
0.827586206896552	0.0344827586206897	0.178526301432776\\
0.827586206896552	0.0689655172413793	0.102487425681528\\
0.827586206896552	0.103448275862069	0.0549843463656875\\
0.827586206896552	0.137931034482759	0.0478539989962521\\
0.827586206896552	0.172413793103448	0.0648798209546963\\
0.827586206896552	0.206896551724138	0.0727055276817409\\
0.827586206896552	0.241379310344828	0.0593752814623226\\
0.827586206896552	0.275862068965517	0.0482718844806363\\
0.827586206896552	0.310344827586207	0.0662467748762002\\
0.827586206896552	0.344827586206897	0.127780534044632\\
0.827586206896552	0.379310344827586	0.195702753587634\\
0.827586206896552	0.413793103448276	0.226964268338378\\
0.827586206896552	0.448275862068965	0.224027795742096\\
0.827586206896552	0.482758620689655	0.203799096187405\\
0.827586206896552	0.517241379310345	0.175019894705015\\
0.827586206896552	0.551724137931034	0.128374699193501\\
0.827586206896552	0.586206896551724	0.0862754687683904\\
0.827586206896552	0.620689655172414	0.0717968556664077\\
0.827586206896552	0.655172413793103	0.0780843146572629\\
0.827586206896552	0.689655172413793	0.0961696523701596\\
0.827586206896552	0.724137931034483	0.124879246463413\\
0.827586206896552	0.758620689655172	0.16784291333739\\
0.827586206896552	0.793103448275862	0.202289935785254\\
0.827586206896552	0.827586206896552	0.213471275262316\\
0.827586206896552	0.862068965517241	0.209911699349592\\
0.827586206896552	0.896551724137931	0.191905412128232\\
0.827586206896552	0.931034482758621	0.159350614421998\\
0.827586206896552	0.96551724137931	0.112216839515538\\
0.827586206896552	1	0.0504437971928324\\
0.862068965517241	0	0.27288775906587\\
0.862068965517241	0.0344827586206897	0.221160941252218\\
0.862068965517241	0.0689655172413793	0.174865038824343\\
0.862068965517241	0.103448275862069	0.133934690477768\\
0.862068965517241	0.137931034482759	0.091862692679063\\
0.862068965517241	0.172413793103448	0.0577670574270774\\
0.862068965517241	0.206896551724138	0.0506083661578154\\
0.862068965517241	0.241379310344828	0.0829611369054746\\
0.862068965517241	0.275862068965517	0.131935938736479\\
0.862068965517241	0.310344827586207	0.168012086110319\\
0.862068965517241	0.344827586206897	0.193001216953702\\
0.862068965517241	0.379310344827586	0.211100672117063\\
0.862068965517241	0.413793103448276	0.220606615699743\\
0.862068965517241	0.448275862068965	0.222183510960154\\
0.862068965517241	0.482758620689655	0.21490728992174\\
0.862068965517241	0.517241379310345	0.195771516708828\\
0.862068965517241	0.551724137931034	0.161212466658194\\
0.862068965517241	0.586206896551724	0.126009156266663\\
0.862068965517241	0.620689655172414	0.101404279590182\\
0.862068965517241	0.655172413793103	0.0801174451629867\\
0.862068965517241	0.689655172413793	0.0762441590252685\\
0.862068965517241	0.724137931034483	0.110813320819548\\
0.862068965517241	0.758620689655172	0.180386316853105\\
0.862068965517241	0.793103448275862	0.235822529413014\\
0.862068965517241	0.827586206896552	0.241299565890537\\
0.862068965517241	0.862068965517241	0.216417330298784\\
0.862068965517241	0.896551724137931	0.183012277828506\\
0.862068965517241	0.931034482758621	0.151525721580986\\
0.862068965517241	0.96551724137931	0.114447291035736\\
0.862068965517241	1	0.0716530339650465\\
0.896551724137931	0	0.243537380759883\\
0.896551724137931	0.0344827586206897	0.246909213110076\\
0.896551724137931	0.0689655172413793	0.234082718322021\\
0.896551724137931	0.103448275862069	0.204701624227627\\
0.896551724137931	0.137931034482759	0.136893983695264\\
0.896551724137931	0.172413793103448	0.0609733340188498\\
0.896551724137931	0.206896551724138	0.0395901443051213\\
0.896551724137931	0.241379310344828	0.104057510665703\\
0.896551724137931	0.275862068965517	0.195790492887219\\
0.896551724137931	0.310344827586207	0.24295030663853\\
0.896551724137931	0.344827586206897	0.23710774896093\\
0.896551724137931	0.379310344827586	0.215033385617524\\
0.896551724137931	0.413793103448276	0.208243557675488\\
0.896551724137931	0.448275862068965	0.216583549081979\\
0.896551724137931	0.482758620689655	0.223201135499048\\
0.896551724137931	0.517241379310345	0.213864297867813\\
0.896551724137931	0.551724137931034	0.190063498865594\\
0.896551724137931	0.586206896551724	0.160752842459076\\
0.896551724137931	0.620689655172414	0.127780783324445\\
0.896551724137931	0.655172413793103	0.0835061504057081\\
0.896551724137931	0.689655172413793	0.0604200418818417\\
0.896551724137931	0.724137931034483	0.0968009883302891\\
0.896551724137931	0.758620689655172	0.183464975036525\\
0.896551724137931	0.793103448275862	0.252685437518866\\
0.896551724137931	0.827586206896552	0.253629829457051\\
0.896551724137931	0.862068965517241	0.213748723097178\\
0.896551724137931	0.896551724137931	0.171006585210607\\
0.896551724137931	0.931034482758621	0.14375198847715\\
0.896551724137931	0.96551724137931	0.118873395658598\\
0.896551724137931	1	0.09620427233353\\
0.931034482758621	0	0.203336744923941\\
0.931034482758621	0.0344827586206897	0.241614196048353\\
0.931034482758621	0.0689655172413793	0.252883481596779\\
0.931034482758621	0.103448275862069	0.236681752122832\\
0.931034482758621	0.137931034482759	0.166365494051419\\
0.931034482758621	0.172413793103448	0.078790894245306\\
0.931034482758621	0.206896551724138	0.0499471145138961\\
0.931034482758621	0.241379310344828	0.113675052739907\\
0.931034482758621	0.275862068965517	0.205883493527537\\
0.931034482758621	0.310344827586207	0.250602807163901\\
0.931034482758621	0.344827586206897	0.237617163963397\\
0.931034482758621	0.379310344827586	0.208518061239641\\
0.931034482758621	0.413793103448276	0.20043153338762\\
0.931034482758621	0.448275862068965	0.213737851128365\\
0.931034482758621	0.482758620689655	0.227175403737476\\
0.931034482758621	0.517241379310345	0.222947087196334\\
0.931034482758621	0.551724137931034	0.204422379492983\\
0.931034482758621	0.586206896551724	0.177773216050444\\
0.931034482758621	0.620689655172414	0.140808306095226\\
0.931034482758621	0.655172413793103	0.0858977747224461\\
0.931034482758621	0.689655172413793	0.0526434508500727\\
0.931034482758621	0.724137931034483	0.0851696054630077\\
0.931034482758621	0.758620689655172	0.171902495858163\\
0.931034482758621	0.793103448275862	0.24214427153876\\
0.931034482758621	0.827586206896552	0.242704611519604\\
0.931034482758621	0.862068965517241	0.202257836230388\\
0.931034482758621	0.896551724137931	0.161658405845283\\
0.931034482758621	0.931034482758621	0.140680187779474\\
0.931034482758621	0.96551724137931	0.125214078204007\\
0.931034482758621	1	0.115093014139233\\
0.96551724137931	0	0.146147566010704\\
0.96551724137931	0.0344827586206897	0.215410899223392\\
0.96551724137931	0.0689655172413793	0.250788740711741\\
0.96551724137931	0.103448275862069	0.25179588315227\\
0.96551724137931	0.137931034482759	0.192157138739506\\
0.96551724137931	0.172413793103448	0.108148895845777\\
0.96551724137931	0.206896551724138	0.0743075011737961\\
0.96551724137931	0.241379310344828	0.118252658652857\\
0.96551724137931	0.275862068965517	0.186530002188312\\
0.96551724137931	0.310344827586207	0.219943011088686\\
0.96551724137931	0.344827586206897	0.21062733487588\\
0.96551724137931	0.379310344827586	0.190820931229676\\
0.96551724137931	0.413793103448276	0.189605075425377\\
0.96551724137931	0.448275862068965	0.208980880758936\\
0.96551724137931	0.482758620689655	0.227907117917728\\
0.96551724137931	0.517241379310345	0.227568610823424\\
0.96551724137931	0.551724137931034	0.211813335051478\\
0.96551724137931	0.586206896551724	0.18619010416703\\
0.96551724137931	0.620689655172414	0.14773392664973\\
0.96551724137931	0.655172413793103	0.0889783879420746\\
0.96551724137931	0.689655172413793	0.0500893134597548\\
0.96551724137931	0.724137931034483	0.0742520816332069\\
0.96551724137931	0.758620689655172	0.149403300685114\\
0.96551724137931	0.793103448275862	0.211882005953649\\
0.96551724137931	0.827586206896552	0.214074673523644\\
0.96551724137931	0.862068965517241	0.181688560769083\\
0.96551724137931	0.896551724137931	0.15083270132075\\
0.96551724137931	0.931034482758621	0.138978489682109\\
0.96551724137931	0.96551724137931	0.13367157897848\\
0.96551724137931	1	0.134771410329804\\
1	0	0.0718844923646628\\
1	0.0344827586206897	0.168430332111424\\
1	0.0689655172413793	0.22805475138922\\
1	0.103448275862069	0.250333088684277\\
1	0.137931034482759	0.214426635508374\\
1	0.172413793103448	0.149008602982936\\
1	0.206896551724138	0.112575684062759\\
1	0.241379310344828	0.117875211445141\\
1	0.275862068965517	0.138048564498595\\
1	0.310344827586207	0.15134981889542\\
1	0.344827586206897	0.156347436484619\\
1	0.379310344827586	0.161929763619637\\
1	0.413793103448276	0.175662414176136\\
1	0.448275862068965	0.202249841165244\\
1	0.482758620689655	0.225409907255164\\
1	0.517241379310345	0.22778859530727\\
1	0.551724137931034	0.212335259508091\\
1	0.586206896551724	0.186123338871798\\
1	0.620689655172414	0.148653039795373\\
1	0.655172413793103	0.0927706683494239\\
1	0.689655172413793	0.0527211221621809\\
1	0.724137931034483	0.0640264294477566\\
1	0.758620689655172	0.116014612139097\\
1	0.793103448275862	0.161997072867216\\
1	0.827586206896552	0.167810384966086\\
1	0.862068965517241	0.152035557045455\\
1	0.896551724137931	0.138473947015568\\
1	0.931034482758621	0.138602680468048\\
1	0.96551724137931	0.14424900471542\\
1	1	0.155325701152984\\
};
\addplot3 [color=red,solid,line width=2.0pt]
 table[row sep=crcr] {0.1	0.5	3\\
0.126027713543254	0.520804667247123	3\\
0.151954449779096	0.538427592773893	3\\
0.180956934635466	0.554444075491578	3\\
0.217149432112214	0.571122948315596	3\\
0.264740450469279	0.58934629452416	3\\
0.303398547255194	0.592224496030499	3\\
0.334576338751455	0.579183426645179	3\\
0.365642932273283	0.556252580040813	3\\
0.396356278738597	0.527018280225036	3\\
0.421867478527032	0.502033996517147	3\\
0.445986233549475	0.480378361434496	3\\
0.467945868580747	0.455464928358652	3\\
0.486100604245849	0.424907475834193	3\\
0.495566814811417	0.390080843549739	3\\
0.495585109204872	0.352984113272876	3\\
0.504399809826646	0.382476086240106	3\\
0.521153560933272	0.406549245977612	3\\
0.544143296348686	0.42150561230869	3\\
0.567748066307488	0.433737657282683	3\\
0.592292766672752	0.446024742845228	3\\
0.619052304341464	0.460025147971715	3\\
0.651649535911148	0.476707599114684	3\\
0.688696457331808	0.49060718060866	3\\
0.722550029959221	0.492995172263751	3\\
0.755743520934362	0.487310994945957	3\\
0.790990571161553	0.47876192948511	3\\
0.827635518907373	0.47738380118364	3\\
0.86405021384148	0.485533193735811	3\\
0.9	0.5	3\\
};
 \end{axis}
\end{tikzpicture}%
\end{document}
%\input{images/path3.tex}
\caption{Optimized flight path from the coordinates 0.1,0.5 to 0.9,0.5, in the time invariant (top) and the time variant (bottom) case.}
\label{fig:path}
\end{figure}
The optimized path (shown in red) for the time invariant case is shown in figure~\ref{fig:path}. We observe that the path leads nicely trough sunny weather grids including a detour to the sunniest spot on the grid at (0.5,0.35), indicating that our problem is well formulated. The optimized solution for the time variant case is shown in the bottom of figure~\ref{fig:path}. Now the path stays longer in the sunny area in the middle and the plane does not take the detour to the center. However some not optimal seeming detours appear. It is important hoever to keep in mind, that this problem is not convex... 









\end{document}
